\documentclass{report}
\usepackage[utf8]{inputenc}
\usepackage[T1]{fontenc}
\usepackage{geometry}
\usepackage{graphicx}
\usepackage{float}
\usepackage{hyperref}
\usepackage[english,indonesian]{babel}

\title{Optimasi Pemrosesan Template LaTeX Menggunakan Pendekatan Berbasis AI}
\author{Academic Support Agent}
\date{\today}

\begin{document}
\maketitle

\begin{abstract}

\end{abstract}

\chapter{PENDAHULUAN}


\chapter{TINJAUAN PUSTAKA}


% Original content (commented out):
% % Modelo de Relatório Intercalar do Trabalho Final de Mestrado 
% Mestrado em Engenharia de Eletrónica e Telecomunicações
% Instituto Superior de Engenharia de Lisboa
%
% V.1.0 2023-03-03 Rui Duarte

\documentclass[a4paper,twoside,11pt]{report}

% Language setting: Replace `english' with `portuguese' to change the document language
%\usepackage[english]{babel}
\usepackage[portuguese]{babel}


% Set page size and margins
% Replace `letterpaper' with `a4paper' for UK/EU standard size
\usepackage[a4paper,top=2cm,bottom=2cm,left=3cm,right=3cm,marginparwidth=1.75cm]{geometry}

% Useful packages
\usepackage{amsmath}
\usepackage{graphicx}
\usepackage[colorlinks=true, linkcolor=black, urlcolor=blue, citecolor=black]{hyperref}
\usepackage{fancyhdr}
\usepackage{acronym}
\usepackage{url}
\usepackage[default]{opensans}
\usepackage{indentfirst}
\setlength{\parindent}{1cm}
\usepackage{setspace}
\usepackage[avantgarde,nogrey]{quotchap}
\usepackage{xcolor}
\colorlet{chaptergrey}{black}
\makeatletter
\renewcommand*{\sectfont}{\bfseries}
\renewcommand*{\chapnumfont}{ \usefont{T1}{\@defaultcnfont}{b}{n}\fontsize{40}{50}\selectfont% Default: 100/130
%  \color{chaptergrey}%
}
\makeatother
\fontsize{11}{13.5}


%************************************************** 

%redefinir a dupla página branca - para não aparecer o cabeçalho sem nada em baixo
\makeatletter
\def\cleardoublepage{\clearpage\if@twoside \ifodd\c@page\else
    \hbox{}
    \vspace*{\fill}
%    \begin{center}
%        This page was intentionally left blank
%    \end{center}
    \vspace{\fill}
    \thispagestyle{empty}
    \newpage
    \if@twocolumn\hbox{}\newpage\fi\fi\fi}
\makeatother

\usepackage{tocloft}
\renewcommand{\cftchapleader}{\cftdotfill{\cftdotsep}} % ... for chapters

\addto\captionsportuguese{\renewcommand{\contentsname}{\hfill Índice}}
\addto\captionsportuguese{\renewcommand{\listfigurename}{\hfill Lista de Figuras}}
\addto\captionsportuguese{\renewcommand{\listtablename}{\hfill Lista de Tabelas}}
\addto\captionsenglish{\renewcommand{\contentsname}{\hfill Table of Contents}}
\addto\captionsenglish{\renewcommand{\listfigurename}{\hfill List of Figures}}
\addto\captionsenglish{\renewcommand{\listtablename}{\hfill List of Tables}}




\begin{document}


\begin{titlepage}
   \begin{center}
       \vspace*{2cm}

    \includegraphics[scale=0.4]{logo.png}\\
       \textbf{INSTITUTO SUPERIOR DE ENGENHARIA DE LISBOA}\\
       \textbf{Departamento de Engenharia Eletrónica e Telecomunicações e Computadores}\\

       \vspace{2cm}
    \includegraphics[width=6cm]{old-phone.jpg}\\
       \vspace{2cm}
       
       {\LARGE \textbf{Título da Dissertação, do Projecto ou do Relatório de Estágio} }
            
       \vspace{1cm}

{ \LARGE{NOME DO CANDIDATO}}\\ %(TNR 16, MAÍUSCULAS, Centrado)
{ \Large{(Grau do candidato)}}\\ %TNR 14p centrado

\vspace{1cm}

Relatório Intercalar do Trabalho Final de Mestrado \\
Mestrado em Engenharia de Eletrónica e Telecomunicações
       \vfill
            
\begin{tabular}{cc}
{Orientador(es): } & \\ %  (TNR 12 p)
& {Grau e Nome }\\ % ((TNR 11 p))
& {Grau e Nome }\\ % (TNR 11 p)
\end{tabular}     

\vspace{1cm}
       
{ \small Mês Ano }
            
   \end{center}
\end{titlepage}

%\cleardoublepage % para impressão em dupla face/livro é necessário deixar a página seguinte vazia 
\clearpage




\doublespacing
\chapter*{Resumo}
Para este relatório intercalar a estrutura do modelo sugerida é idêntica à do documento final do trabalho final de mestrado, embora bastante mais simples. Algumas partes formais relacionadas com agradecimentos, listas de acrónimos, listas de variáveis, lista de símbolos, etc. foram suprimidas para este relatório. O aluno, caso considere necessário poderá neste relatório utilizar algum destes recursos.
No relatório intercalar basta apresentar o resumo em Português. No trabalho final será necessário um resumo na língua do texto principal e um resumo noutra língua. Assume-se que as duas línguas em questão serão sempre o Português e o Inglês.
O resumo não deve exceder uma página e de uma forma genérica deverá responder às seguintes questões:
\begin{enumerate}
    \item Qual é o problema a resolver?
\item	Porque é que é esse problema é interessante / desafiante / pertinente?
\item	Qual é qual a metodologia/solução proposta para o resolver?
\item	O que resulta (implicações/consequências) da solução?
\end{enumerate}

\vspace{1cm}
\noindent \textbf{Palavas chave}: Palavra-chave 1, Palavra-chave 2, Palavra-chave 3, ….


%\cleardoublepage
\clearpage

\doublespacing
\tableofcontents
%\cleardoublepage
\clearpage
\listoffigures
%\cleardoublepage
\clearpage
\listoftables
%\cleardoublepage
\clearpage


\onehalfspacing



\chapter{NOME DO 1º CAPÍTULO}
É pretendido que o relatório intercalar seja um documento curto e sintético, que demonstre a capacidade do aluno de explicar o problema a resolver, a metodologia OU solução para o resolver, o estado da arte/trabalho relacionado e indique ainda uma lista de referências/bibliografia adequada. Pretende-se também que alguns aspetos formais sejam observados. 
A extensão deste documento não deverá exceder 25 páginas.


\section{Organização do documento}
Para estruturar melhor o documento poderá usar subsecções e eventualmente ainda sub-subsecções. Em cada subsecção deverá apresentar-se detalhes do tema da secção a que pertence, caso se considere necessário organizar o texto dessa forma. Não é de todo obrigatório fazê-lo, nem recomendável se a informação do capítulo não o justificar, correndo-se o risco de ter um parágrafo em cada subsecção ou sub-subsecção, caso a informação seja reduzida.
Os capítulos devem iniciar-se sempre em páginas ímpares.

\subsection{Estrutura do relatório intercalar}
Para o relatório intercalar o documento a apresentar deverá conter:
\begin{enumerate}
    \item Capa (Versão simplificada da capa proposta pelo ISEL sem indicação do júri)
    \item 	Resumo e palavras-chave
    \item 	Índice
    \item 	Lista de figuras (se houver mais de três)
    \item 	Lista de tabelas (se houver mais de três)
    \item 	Texto principal organizado em capítulos 
    \item 	Bibliografia
\end{enumerate}

\subsection{Figuras}
As figuras do documento deverão ser referidas no corpo do texto antes de serem apresentadas. É conveniente usar indexação automática para assegurar que a referenciação de figuras está correta. Todas as figuras devem ter uma legenda. A referência bibliográfica correspondente à origem da figura deve ser mencionada. A título de exemplo, para este caso, uma possível descrição antes de inserir a figura seria: “Na Figura \ref{fig:my_label1} mostra-se uma antena de...”.

\begin{figure}[tbp]
    \centering
    \includegraphics[scale=0.3]{fig1.png}
    \caption{Exemplo de figura e da legenda}
    \label{fig:my_label1}
\end{figure} 

Depois da apresentação da figura, deverá seguir-se o comentário sobre o seu conteúdo por forma a contribuir para a narrativa a ela associada. A legenda da figura deverá estar centrada na linha, a menos que seja muito longa e ocupe mais que uma linha. Nesse caso, a legenda deverá ficar justificada em vez de centrada. O mesmo se aplica às legendas das tabelas.

\subsection{Tabelas}
As tabelas obedecem a uma lógica idêntica à das figuras, com numeração, legenda e mencionadas antes de serem apresentadas.
Tabela \ref{tab:my_label} — Exemplo de tabela e da sua legenda.


\begin{table}[tbp]
    \centering
    \caption{Exemplo de tabela e da sua legenda.}
    \begin{tabular}{|c|c|}
        \hline
        \textbf{Material} &	\textbf{Índice de refração} \\
        \hline
        Ar       & 	1,0003 \\
        água     &	1,33 \\
        Gelo     &	1,31 \\
        Álcool   &	1,36 \\
        Vidro Crown &	1,52 \\
        Diamante    &	2,42 \\
        \hline
    \end{tabular}
    \label{tab:my_label}
\end{table}


\section{Referenciação Cruzada}
Sempre que for necessário referir itens numerados do texto (figuras, tabelas, equações, etc.), deve-se inserir uma referência cruzada usando o menu de indexação automática. Assim, se os itens forem renumerados (por exemplo, porque foi inserida uma nova figura no meio de duas figuras existentes), as referências também serão atualizadas automaticamente.







%\cleardoublepage
\clearpage





\chapter{NOME DO 2ºCAPÍTULO}

\section{Numeração das Figuras}
Recomenda-se a numeração das figuras associadas ao capítulo a que pertencem, conforme se mostra na Figura \ref{fig:my_label2}.

 \begin{figure}[tbp]
     \centering
     \includegraphics[scale=0.7]{fig2.png}
     \caption{Exemplo de figura noutro capítulo e da legenda.}
     \label{fig:my_label2}
 \end{figure}
 

\section{Bibliografia}
A bibliografia deve ser apresentada no final do documento, a seguir ao corpo do texto do último capítulo, e antes dos anexos que possam existir. Esta secção do documento deve iniciar-se numa nova página ímpar. Será conveniente que no documento final do Trabalho Final de Mestrado, a bibliografia esteja coerente com as regras e modelos de citação que a seguir se descrevem, podendo, no relatório intercalar, estes aspetos formais não serem seguidos com tanto rigor.


\subsection{Estilos Bibliográficos}
Existem muitos padrões e estilos bibliográficos. Cada área científica tem sua própria forma de apresentar tanto as citações quanto as referências bibliográficas. Os estilos mais comuns são o APA (American Psychological Association - autor/data), 7ª edição, e o IEEE (Institute of Electrical and Electronics Engineers - numérico). 

\subsection{Modelo de citação}
Existem várias formas de citar outros autores, podendo-se estas dividir-se em duas grandes classes:
\begin{itemize}
    \item Citações indiretas ou conceituais, nas quais reproduzimos as ideias de outra pessoa pelas nossas próprias palavas;
\item	Citações diretas ou formais, nas quais transcrevemos exatamente as palavras de um autor usando aspas.
Os modelos de citação seguem 3 sistemas:
\item Sistema autor-data, em que a citação aparece assim: (Santos, 2003), se forem dois autores (Santos e Correia, 2003) e se forem mais de 5 autores (Santos, et al., 2003), sendo o mais conhecido e utilizado o estilo APA;
\item Sistema numérico, em que cada citação é identificada com um número [1] e a lista de referências bibliográficas é compilada no final do trabalho (bibliografia), sendo o estilo mais conhecido e utilizado o IEEE.
\item Sistemas mistos, nos quais a citação no texto é numérica, mas a bibliografia é ordenada alfabeticamente pelo sobrenome do autor. Exemplos de estilos mistos são: Springer Lecture notes in Computer Science (ordenado alfabeticamente) e o Council of Science Editors, Citation-Name (ordenado alfabeticamente numérico), entre outros.
\end{itemize}

Os modelos de citação mais utilizados, em geral, são APA e IEEE. Recomenda-se a utilização do modelo IEEE, a menos que o orientador tenha outra recomendação.

\subsection{Modelo IEEE}
O modelo de citação \href{https://ieee-dataport.org/sites/default/files/analysis/27/IEEE\%20Citation\%20Guidelines.pdf}{IEEE} 
inclui regras para vários tipos de documentos. 
As referências seguintes ilustram o modelo IEEE para algumas das fontes de citação mais utilizadas, apresentando-se o formato genérico e um exemplo de cada caso: 

\begin{itemize}
    \item livro \cite{r1} e exemplo de livro \cite{r2}; 
\item 	livro online e exemplo \cite{r3},\cite{r4}; 
\item 	capítulo de livro \cite{r5},\cite{r6}; 
\item 	artigo científico \cite{r7},\cite{r8},\cite{r9};
\item 	artigo em atas de conferência \cite{r10},\cite{r11}, 
\item 	website \cite{r12}; 
\item 	software \cite{r13},\cite{r14}.
\end{itemize}


Mais fontes (manuais, aulas, vídeos, etc.) poderão ser consultadas nas recomendações IEEE.










%\cleardoublepage
\clearpage


\bibliographystyle{ieee}
%\bibliography{sample} % carregar a partir de um ficheiro .bib 
% recomenda-se uso de um gestor de referencias bibliograficas tal como:
% Mendeley ou Jabref
% a alternativa passa por listar todas as referências com \bibitem{cite_id}Citação
\begin{thebibliography}{9}
\bibitem{r1}J. K. Author, “Title of chapter in the book,” in Title of His Published Book, xth ed. City of Publisher, (only U.S. State), Country: Abbrev. of Publisher, year, ch. x, sec. x, pp. xxx–xxx.

\bibitem{r2}R. L. Myer, “Parametric oscillators and nonlinear materials,” in Nonlinear Optics, vol. 4, P. G. Harper and B. S. Wherret, Eds. San Francisco, CA, USA: Academic, 1977, pp. 47–160.

\bibitem{r3}J. K. Author, “Title of chapter in the book,” in Title of Published Book, xth ed. City of Publisher, State, Country: Abbrev. of Publisher, year, ch. x, sec. x, pp. xxx–xxx. [Online]. Available: \url{http://www.web.com}.

\bibitem{r4}G. O. Young, “Synthetic structure of industrial plastics,” in Plastics, vol. 3, Polymers of Hexadromicon, J. Peters, Ed., 2nd ed. New York, NY, USA: McGraw-Hill, 1964, pp. 15–64. [Online]. Available: \url{http://www.bookref.com}.

\bibitem{r5}T. Ogura, “Electronic government and surveillance-oriented society,” in Theorizing Surveillance: The Panopticon and Beyond. Cullompton, U.K.: Willan, 2006, ch. 13, pp. 270–295.

\bibitem{r6}J. K. Author, “Title of chapter in the book,” in Title of His Published Book, xth ed. City of Publisher, (only U.S. State), Country: Abbrev. of Publisher, year, ch. x, sec. x, pp. xxx–xxx.

\bibitem{r7}J. K. Author, “Name of paper,” Abbrev. Title of Periodical, vol. x, no. x, pp. xxx–xxx, Abbrev. month, year, doi: xxx. 

\bibitem{r8}M. M. Chiampi and L. L. Zilberti, “Induction of electric field in human bodies moving near MRI: An efficient BEM computational procedure,” IEEE Trans. Biomed. Eng., vol. 58, no. 10, pp. 2787–2793, Oct. 2011, doi: 10.1109/TBME.2011.2158315.

\bibitem{r9}K. A. Nelson, R. J. Davis, D. R. Lutz, and W. Smith, “Optical generation of tunable ultrasonic waves,” Journal of Applied Physics, vol. 53, no. 2, Feb., pp. 1144-1149, 2002.

\bibitem{r10}J. K. Author, “Title of paper,” in Abbreviated Name of Conf., (location of conference is optional), year, pp. xxx– xxx, doi: xxx. 

\bibitem{r11}G. Veruggio, “The EURON roboethics roadmap,” in Proc. Humanoids ’06: 6th IEEE-RAS Int. Conf. Humanoid Robots, 2006, pp. 612–617, doi: 10.1109/ICHR.2006.321337.


\bibitem{r12}European Telecommunications Standards Institute, “Digital Video Broadcasting (DVB): Implementation guide for DVB terrestrial services; transmission aspects,” European Telecommunications Standards Institute, ETSI-TR-101, 2007. [Online]. Available: \url{http://www.etsi.org}. [Accessed: Nov. 12, 2007].

\bibitem{r13}Title of Software. (version or year), Publisher Name. Accessed: Date (when applicable). [Type of Medium]. Available: \url{site/path/file}

\bibitem{r14}Ngspice. (2011). [Online]. Available: \url{http://ngspice.sourceforge.net}

\end{thebibliography}



\end{document}

\end{document}